% !TeX root = ../../Book1.tex
% 纲要标题：### 4.3 轨道计数
% 纲要位置：计划纲要.md，第 229 行。

\section{轨道计数}
\label{b1:sec:0403}

当若干对象在对称变换下视为相同，所求数量就是轨道数。不能一概用对象总数除以群阶，因为对称性较强的对象有较小的轨道。直接列出每个轨道往往困难，固定一个变换后计算它保持不动的对象却容易得多。对这些不动点数取平均，将给出正确的计数公式。

% 本节正文按下列原纲要要点撰写。

% 纲要内容（仅作写作计划，尚非教材正文）：
%
% * Burnside 轨道计数引理
% * 染色与有限组合计数
% * 作用的核与忠实性
%

\subsection{Burnside 轨道计数引理}
\label{b1:subsec:0403:01}

下面的\term[burnside-guidao-jishu-yinli]{Burnside 轨道计数引理}[Burnside orbit-counting lemma]把轨道数与各个群元素的不动点数联系起来。

\begin{theorem}[Burnside 轨道计数]\label{b1:thm:burnside-count}
设有限群 \(G\) 作用在有限集合 \(X\) 上。记轨道集合为 \(X/G\)，并记
\[
X^g\coloneqq\{\,x\in X\mid gx=x\,\}.
\]
则
\[
|X/G|=\frac{1}{|G|}\sum_{g\in G}|X^g|.
\]
\symidx{\(X^g\)：一个群元素的不动点集}
\symidx{\(X/G\)：作用的轨道集合}
\end{theorem}
\begin{proof}
考察集合 \(E=\bigl\{\,(g,x)\in G\times X\mid gx=x\,\bigr\}\)。先固定 \(g\)，得到
\(|E|=\sum_g|X^g|\)；先固定 \(x\)，得到
\(|E|=\sum_x|G_x|\)。
对同一轨道 \(O\) 中的每个 \(x\)，\cref{b1:thm:orbit-stabilizer}给出 \(|G_x|=|G|/|Gx|=|G|/|O|\)，所以该轨道对后一和式的贡献为
\[
\sum_{x\in O}|G_x|=|O|\frac{|G|}{|O|}=|G|.
\]
每个轨道贡献一次 \(|G|\)，故 \(|E|=|X/G|\,|G|\)，除以 \(|G|\) 即得公式。
\end{proof}
公式取的是不动点数的平均值；即使某些轨道较小，也不需要另外为它们加权。直接把对象总数除以群阶只有在所有稳定子均平凡时才正确。一个完全对称的对象可能被很多变换固定，双重计数正好补偿这种轨道大小的差异。

\subsection{染色与有限组合计数}
\label{b1:subsec:0403:02}

给正方形四个顶点涂上 \(k\) 种有名称的颜色，允许重复使用颜色。把顶点集合记为 \(V\)，颜色集合记为 \(C\)，全部染色组成 \(X=C^V\)，即从 \(V\) 到 \(C\) 的函数集合。若 \(G\) 是所允许的对称群，正确的左作用是
\[
(g\cdot f)(v)=f(g^{-1}v).
\]
逆元使复合方向相容：\(g\cdot(h\cdot f)\) 在 \(v\) 的值为
\(f(h^{-1}g^{-1}v)=f\bigl((gh)^{-1}v\bigr)\)。染色 \(f\) 被 \(g\) 固定，当且仅当它在 \(g\) 的每个顶点轮换上取常值。因此，若 \(g\) 在 \(V\) 上有 \(c(g)\) 个轮换（包括不动顶点），则 \(|X^g|=k^{c(g)}\)。

\ProseParagraphBreak
若只允许旋转，群为 \(C_4\)。单位有四个轮换，两个四分之一周旋转各有一个轮换，半周旋转有两个轮换。把相应的不动点数 \(k^4,k,k,k^2\) 代入\cref{b1:thm:burnside-count}并除以群阶 \(4\)，得到染色轨道数
\[
\frac{k^4+k^2+2k}{4}.
\]
例如 \(k=2\) 时得到 \(6\) 类。颜色名称不随对称变换交换，所以全红与全蓝是两个不同的轨道。

\ProseParagraphBreak
若允许全部正方形对称，群为 \(D_8\)。两条对角线反射各有两个固定顶点和一个二轮换，因此有三个轮换；两条对边中点连线反射各有两个二轮换。在四个旋转的不动点总数之外再加上 \(2k^3+2k^2\)，\cref{b1:thm:burnside-count}要求这次除以 \(|D_8|=8\)，所以计数变为
\[
\frac{k^4+2k^3+3k^2+2k}{8}.
\]
当 \(k=3\) 时，旋转等价下有 \(24\) 类，全部对称等价下有 \(21\) 类。两种答案的不同来自等价关系不同，不是同一个计数问题的矛盾答案。

\ProseParagraphBreak
有额外限制时，应先把 \(X\) 换成满足限制的染色集合，再检查该集合是否被群保持。例如要求相邻顶点异色，这个条件在 \(D_8\) 下保持，但此时不能直接使用 \(k^{c(g)}\)：有的轮换会迫使相邻顶点同色，固定染色数便可能为零。轨道计数引理仍成立，需要修改的是逐个变换的不动点计算。

\subsection{作用的核与忠实性}
\label{b1:subsec:0403:03}

\begin{definition}\label{b1:def:faithful-action}
设群 \(G\) 作用在集合 \(X\) 上。使 \(X\) 中每一点都不动的群元素组成作用的核，记为
\[
K\coloneqq\{\,g\in G\mid gx=x\text{ 对所有 }x\in X\,\}.
\]
若只有单位元素固定全部的点，即 \(K=\{1\}\)，称此作用为\term[zhongshi-zuoyong]{忠实作用}[faithful action]。若对任意 \(x\in X\) 和 \(g\in G\)，等式 \(gx=x\) 都蕴含 \(g=1\)，称此作用为\term[ziyou-zuoyong]{自由作用}[free action]。
\end{definition}
由定义，\(K\) 等于对应同态 \(\rho:G\rightarrow S_X\) 的核，也等于所有稳定子的交 \(\bigcap_{x\in X}G_x\)。\cref{b1:prop:kernel-fibers}说明同态单射当且仅当核平凡，因此忠实等价于 \(\rho\) 单射；自由则按定义等价于每个稳定子都平凡。

\ProseParagraphBreak
自由作用一定忠实（这里假设 \(X\neq\varnothing\)），反过来不成立：\(S_3\) 在三个点上的自然作用忠实，但换位仍固定一个点。忠实性要求一个元素只要固定所有点就必须为单位；自由性要求一个元素只要固定任一点就必须为单位，二者的量词不同。

\begin{proposition}\label{b1:prop:quotient-action}
若 \(K\) 是作用的核，则公式 \((gK)x=gx\) 定义 \(G/K\) 在 \(X\) 上的忠实作用，且与原作用具有相同的轨道。
\end{proposition}
\begin{proof}
首先，\(K=\ker\rho\) 是同态的核，故由\cref{b1:prop:normal-basic-examples}正规，\cref{b1:thm:quotient-group}保证 \(G/K\) 确为群。若 \(gK=hK\)，写 \(g=hk\)、\(k\in K\)，则 \(gx=h(kx)=hx\)，所以定义与代表元无关。两条作用公理从原作用继承。若 \(gK\) 固定每个点，则 \(g\in K\)，所以这个陪集是单位。最后，任一 \(g\) 与其陪集对 \(X\) 的变换相同，故轨道不变。
\end{proof}
例如 \(D_8\) 对两条对角线的作用有核 \(\{1,r^2,s,r^2s\}\)，这里 \(s\) 取沿一条对角线的反射。商群 \(C_2\) 忠实地交换两条对角线。这个两点作用不能区分八个对称，却足以发现一个指数为 \(2\) 的正规子群。Burnside 公式不要求忠实性；若按商群求和，每个商元素的不动点数在原群中恰重复 \(|K|\) 次，平均值保持不变。
